\chapter{Refined Products and Cracks}\label{ch:m3:refined-products-and-cracks}

A refinery on the US Gulf Coast buys crude and sells gasoline and diesel. It does
not care much whether oil costs sixty dollars or a hundred; it cares about the
difference between what its products fetch and what its crude costs. That
difference has a name on every energy screen, the crack, after the cracking of
heavy molecules into light ones. In August 2026 the most common version of it
averaged about \$65 a barrel of crude in New York Harbor, three times its average
of the previous twenty years, and a refiner could have locked most of a quarter's
margin in an afternoon with three futures contracts. This chapter covers the
products, the spreads that price the refining business, the arbitrage that moves
products between regions, and the seasons and regulations that shape both.

\section{The refinery and its product slate}

Distillation separates crude by boiling point: gases and naphtha at the top,
kerosene and gasoil in the middle, heavy fuel oil and residue at the bottom.
Conversion units (catalytic crackers, hydrocrackers, cokers) break the heavy
fractions into lighter, more valuable ones; treating units remove sulphur;
blending combines components into products that meet a specification.

\begin{definition}[Product slate]\label{def:m3:refined-products-and-cracks:slate}
A refinery's \emph{product slate}\index{product slate} is the set of products it
makes and their shares of its output, per barrel of crude processed.
\end{definition}

A US refinery gets about 19 to 20 gallons of gasoline and 11 to 13 gallons of
ultra-low-sulphur distillate from a 42-gallon barrel of crude; the slate varies
with the crude, the season and the units the refinery has. The total output
usually exceeds 42 gallons, because the products are less dense than the crude:
refiners call it processing gain.

\begin{omfigure}
\centering
\begin{tikzpicture}[font=\small,
  box/.style={draw,thick,rounded corners=2pt,align=center,minimum height=0.8cm},
  arr/.style={-{Stealth[length=2.2mm]},thick}]
\node[box,draw=omDef,fill=omDef!8,minimum width=2.2cm] (crude) at (0,0) {crude oil\\42 gallons};
\node[box,draw=omThm,fill=omThm!8,minimum width=1.8cm,minimum height=3.6cm] (col) at (3.6,0) {distillation\\and\\conversion};
\draw[arr] (crude) -- (col);
\node[box,draw=omProp,fill=omProp!8,minimum width=4.4cm,anchor=west] (g) at (6.0,1.35) {gasoline: about 19--20 gallons};
\node[box,draw=omProp,fill=omProp!8,minimum width=4.4cm,anchor=west] (d) at (6.0,0.2) {distillate: about 11--13 gallons};
\node[box,draw=omExo,fill=omExo!8,minimum width=4.4cm,anchor=west] (o) at (6.0,-0.95) {jet fuel, fuel oil, others};
\draw[arr] (col.east |- g) -- (g);
\draw[arr] (col.east |- d) -- (d);
\draw[arr] (col.east |- o) -- (o);
\node[font=\footnotesize,anchor=north] at (3.6,-2.0) {the products together exceed 42 gallons: processing gain};
\end{tikzpicture}

\omcaption{A barrel of crude through a US refinery. Shares vary by crude, season and
refinery configuration. Source of the gasoline and distillate volumes: US Energy
Information Administration. Schematic.}\label{fig:m3:refined-products-and-cracks:slate}
\end{omfigure}

\section{Gasoline, distillates and fuel oil}

\begin{definition}[Middle distillates]\label{def:m3:refined-products-and-cracks:middle}
\emph{Middle distillates}\index{middle distillates} are the products of the middle
of the distillation range: kerosene and jet fuel, diesel and heating oil, and the
gasoils from which they are made.
\end{definition}

Each product has its own benchmarks. In the United States the futures are
gasoline blendstock (RBOB) and ultra-low-sulphur diesel (ULSD), both delivered in
New York Harbor in lots of 42{,}000 gallons (1{,}000 barrels) and quoted in dollars
per gallon; in Europe, low-sulphur gasoil, delivered in the Amsterdam--Rotterdam--Antwerp
area in lots of 100 tonnes and quoted in dollars per tonne. A tonne of gasoil of
density \qty{0.845}{kg/l} is 1.1835 cubic metres, 7.44 barrels: to compare it with crude,
divide the price per tonne by 7.44.

\begin{definition}[Marine fuel sulphur cap]\label{def:m3:refined-products-and-cracks:imo}
The \emph{marine fuel sulphur cap}\index{marine fuel sulphur cap} is the
International Maritime Organization's limit on the sulphur content of fuel used by
ships: 0.50\% by mass since 1 January 2020 outside emission control areas (down
from 3.50\%), and 0.10\% inside them since 2015.
\end{definition}

The cap moved demand from high-sulphur fuel oil, the heavy residue that ships had
burned for decades, towards low-sulphur blends and marine gasoil. It is the
cleanest example of a regulation that re-priced a product overnight: the value of
a refinery's residue depended from one day to the next on whether it had the units
to desulphurise it.

\section{Crack spreads and refining margins}

\begin{definition}[Crack spread, 3-2-1 crack spread]\label{def:m3:refined-products-and-cracks:crack}
A \emph{crack spread}\index{crack spread} is the difference between the value of
refined products and the cost of the crude they are made from, in a fixed volume
ratio, per barrel of crude. The \emph{3-2-1 crack spread}\index{3-2-1 crack spread}
sells two barrels of gasoline and one of diesel for three barrels of crude:
\[
\mathrm{C}_{321} = \frac{2\,P_{\mathrm{gas}} + P_{\mathrm{dist}} - 3\,P_{\mathrm{crude}}}{3},
\]
all prices in dollars per barrel (product prices per gallon times 42).
\end{definition}

\begin{example}[A crack from three quotes]\label{ex:m3:refined-products-and-cracks:crack}
With crude at \qty{83.90}{\$/\barrel}, gasoline at \qty{3.2132}{\$/gal} and diesel at
\qty{4.2530}{\$/gal} (the August 2026 averages of the data used below), the products
are worth \qty{134.95}{\$/\barrel} and \qty{178.63}{\$/\barrel}, and the 3-2-1 crack is
\qty{65.61}{\$/\barrel}. Separately, the gasoline crack is \qty{51.06}{} and the diesel
crack \qty{94.73}{\$/\barrel}.
\end{example}

A crack is a stylised refinery: a real one has its own slate and costs.

\begin{definition}[Gross refining margin]\label{def:m3:refined-products-and-cracks:grm}
The \emph{gross refining margin}\index{gross refining margin} of a refinery is the
value of the products it obtains from a barrel of a given crude, at their market
prices and yields (including processing gain), less the price of that crude. The net
margin also deducts operating costs.
\end{definition}

\begin{proposition}[A crack is a portfolio of futures]\label{prop:m3:refined-products-and-cracks:hedge}
A refiner that processes $Q$ barrels of crude in a period and sells products in the
ratio $2:1$ locks the 3-2-1 crack prevailing today by buying $Q$ barrels of crude
futures and selling $2Q/3$ of gasoline and $Q/3$ of diesel futures for that period.
Its margin then moves only with the difference between its own yields, locations and
grades and those of the contracts.
\end{proposition}

\begin{proof}
On the hedged volume the refiner's physical margin per barrel is $\frac{1}{3}(2P_g + P_d
- 3P_c)$ at the period's prices; the futures gain $\frac{1}{3}(2(F_g - P_g) + (F_d -
P_d) - 3(F_c - P_c))$ per barrel if the physical prices equal the futures at
delivery. The sum is the crack at today's futures prices.
\end{proof}

\begin{omfigure}
\begin{tikzpicture}
\begin{axis}[width=11cm,height=5cm,
  ylabel={\$ per barrel of crude},
  tick label style={font=\small},label style={font=\small},
  xmin=2006.5,xmax=2026.75,ymin=0,ymax=100,grid=major,grid style={gray!25},
  xtick={2008,2010,2012,2014,2016,2018,2020,2022,2024,2026},xticklabel style={/pgf/number format/1000 sep={}},
  legend columns=3,legend style={font=\footnotesize,at={(0.5,-0.18)},anchor=north,draw=none,fill=none,/tikz/every even column/.append style={column sep=8pt}},
  table/col sep=comma]
\addplot[omDef,very thick] table[x=t,y=crack]{figdata/markets-3/03-refined-products-and-cracks/crack321.csv};
\addplot[omProp,thin] table[x=t,y=gas]{figdata/markets-3/03-refined-products-and-cracks/crack321.csv};
\addplot[omThm,thin,dashed] table[x=t,y=diesel]{figdata/markets-3/03-refined-products-and-cracks/crack321.csv};
\legend{3-2-1 crack,gasoline crack,diesel crack}
\end{axis}
\end{tikzpicture}

\omcaption{New York Harbor cracks against WTI, monthly averages of daily spot
prices, July 2006 to August 2026. The peaks are May 2022 (\$62.68) and August 2026
(\$65.61, the highest); the low is November 2009 (\$5.51). Data: FRED series
DCOILWTICO, DGASNYH and DDFUELNYH (US Energy Information
Administration).}\label{fig:m3:refined-products-and-cracks:crack}
\end{omfigure}

The crack averaged \qty{21.20}{\$/\barrel} over the period
(\cref{fig:m3:refined-products-and-cracks:crack}). It rose in 2022, after Russia's full-scale invasion of Ukraine put at risk the source
of half of Europe's diesel imports (the European Union banned seaborne Russian diesel
in early 2023), and again in 2026, when the Gulf disruption of \cref{ch:m3:crude-oil} took refined products as well as crude off
the market; in both episodes diesel cracked far above gasoline.

\begin{dated}{2026-09}{US refining capacity}\label{dat:m3:refined-products-and-cracks:capacity}
US operable atmospheric distillation capacity was 18.2 million barrels per calendar day
on 1 January 2026, down about 250{,}000 barrels a day (about 1\%) from a year earlier,
at 130 operable refineries (301 in 1982).
\end{dated}

\section{Regional arbitrage}

The same product trades at different prices in different places, and moves
between them by pipeline or ship when the difference pays for the move.

\begin{definition}[Arbitrage window]\label{def:m3:refined-products-and-cracks:arb}
An \emph{arbitrage window}\index{arbitrage window} for a product between two regions
is open when its price in the destination exceeds its price at the origin by more
than the cost of moving it (freight, insurance, losses, duties and financing over the
voyage), and closed otherwise.
\end{definition}

The US Gulf Coast refines far more than it consumes and ships the surplus to the
East Coast by pipeline and tanker, and abroad. New York Harbor gasoline therefore
usually trades above Gulf Coast gasoline: by 6.0 cents a gallon on average in the
monthly data since 2006, below it in only 35 of 242 months
(\cref{fig:m3:refined-products-and-cracks:arb}). When the difference falls below the
cost of moving a gallon, the flow slows; when it rises above, more product moves.
\omterm{def:m3:refined-products-and-cracks:arb}{Arbitrage windows} are what connect regional product markets, and what disconnect
them when a pipeline breaks or a canal closes.

\begin{omfigure}
\begin{tikzpicture}
\begin{axis}[width=11cm,height=4.4cm,
  ylabel={cents per gallon},
  tick label style={font=\small},label style={font=\small},
  xmin=2006.5,xmax=2026.75,ymin=-40,ymax=70,grid=major,grid style={gray!25},
  xtick={2008,2010,2012,2014,2016,2018,2020,2022,2024,2026},xticklabel style={/pgf/number format/1000 sep={}},
  table/col sep=comma]
\addplot[omDef,thick] table[x=t,y=cents]{figdata/markets-3/03-refined-products-and-cracks/arb.csv};
\draw[gray] (axis cs:2006.5,0) -- (axis cs:2026.75,0);
\end{axis}
\end{tikzpicture}

\omcaption{New York Harbor minus US Gulf Coast conventional gasoline, monthly averages
of daily spot prices, July 2006 to August 2026. The spread prices the move from the
refining centre to the consuming coast. Data: FRED series DGASNYH and DGASUSGULF (US
Energy Information Administration).}\label{fig:m3:refined-products-and-cracks:arb}
\end{omfigure}

\section{Seasonality and specifications}

Products have seasons. Gasoline demand peaks in the northern summer driving season,
distillate demand in winter heating; refineries schedule maintenance in spring and
autumn between the peaks. The specifications change with the season too.

\begin{definition}[Reid vapour pressure]\label{def:m3:refined-products-and-cracks:rvp}
\emph{Reid vapour pressure}\index{Reid vapour pressure} (RVP) is the standard measure
of gasoline volatility, in pounds per square inch; a lower RVP means less evaporation
and fewer emissions of volatile organic compounds.
\end{definition}

US summer gasoline, sold from 1 June to 15 September, may not exceed an RVP of
\qty{9.0}{psi}, and \qty{7.8}{psi} in some areas. Summer grade needs fewer cheap light
components such as butane and costs more to make, so the gasoline futures contract
changes grade in the spring and its price steps up: the NYMEX rulebook caps the RVP of RBOB delivered at \qty{13.5}{psi} in March and
\qty{7.4}{psi} from April to mid-September.
\Cref{fig:m3:refined-products-and-cracks:season} shows the net effect on the cracks:
from 2010 to 2025 the gasoline crack averaged \qty{2.83}{\$/\barrel} above its mean in
May and \qty{3.83}{\$/\barrel} below it in January; the diesel crack peaked in
November, \qty{4.15}{\$/\barrel} above its mean.

\begin{omfigure}
\begin{tikzpicture}
\begin{axis}[width=11cm,height=4.8cm,ybar,bar width=5pt,
  xlabel={calendar month},ylabel={\$/bbl from the mean},
  tick label style={font=\small},label style={font=\small},
  xmin=0.4,xmax=12.6,ymin=-5,ymax=5,grid=major,grid style={gray!25},
  xtick={1,2,3,4,5,6,7,8,9,10,11,12},
  legend columns=2,legend style={font=\footnotesize,at={(0.5,-0.42)},anchor=north,draw=none,fill=none,/tikz/every even column/.append style={column sep=8pt}},
  table/col sep=comma]
\addplot[fill=omProp!60,draw=omProp] table[x=month,y=gas]{figdata/markets-3/03-refined-products-and-cracks/season.csv};
\addplot[fill=omThm!60,draw=omThm] table[x=month,y=diesel]{figdata/markets-3/03-refined-products-and-cracks/season.csv};
\legend{gasoline crack,diesel crack}
\end{axis}
\end{tikzpicture}

\omcaption{Seasonality of New York Harbor cracks: the mean by calendar month minus the
overall mean, 2010 to 2025. Gasoline is dear in the driving season, diesel in the
heating season. Data: as \cref{fig:m3:refined-products-and-cracks:crack}.}\label{fig:m3:refined-products-and-cracks:season}
\end{omfigure}

\section{Tutorial: cracks from prices, and the hedge}\label{tut:m3:refined-products-and-cracks:crack}

\begin{tutorial}
\textbf{Goal.} Compute cracks from product and crude prices in their own units, and
size the futures hedge of \cref{prop:m3:refined-products-and-cracks:hedge}.
\textbf{End state:} \cref{fig:m3:refined-products-and-cracks:crack,fig:m3:refined-products-and-cracks:season}
and the lots of the weekend problem.
\begin{tutsteps}
\item \textbf{The crack.} Any ratio, with the check that products and crude balance.
\omcode{code/firm/cracks/firm_cracks.py}{27}{38}{A crack spread and the 3-2-1 crack.}\label{lst:m3:refined-products-and-cracks:crack}
\item \textbf{The hedge.} Crude bought, products sold, in whole lots.
\omcode{code/firm/cracks/firm_cracks.py}{47}{54}{Futures lots that lock a crack on a throughput.}\label{lst:m3:refined-products-and-cracks:lots}
\item \textbf{Run} \texttt{m3\_cracks.crack\_series()}, \texttt{seasonality()},
\texttt{refinery\_hedge()} and \texttt{fig\_cracks.py}.
\end{tutsteps}
\textbf{What to change next.} Replace the 3-2-1 by a 5-3-2 and by a European crack on
Brent and gasoil in tonnes; lag the products by a month (the time crude takes to
become product) and see how the crack's volatility changes.
\end{tutorial}

\section{Build: the crack calculator}\label{bld:m3:refined-products-and-cracks:cracks}

\begin{build}
\bfield{Purpose} The miniature firm trades refining margins for refiners' hedges and
for its own book: it needs cracks in any ratio and unit, refinery margins from yields,
hedge sizes in lots, and a test of whether an \omterm{def:m3:refined-products-and-cracks:arb}{arbitrage window} is open.
\bfield{Interface} \texttt{per\_gallon\_to\_per\_barrel}; \texttt{barrels\_per\_tonne(density)};
\texttt{per\_tonne\_to\_per\_barrel}; \texttt{crack(crude, products, ratio, crude\_barrels)};
\texttt{three\_two\_one}; \texttt{gross\_refining\_margin(crude, yields, products)};
\texttt{hedge\_lots(throughput\_bpd, days, ratio, crude\_barrels)};
\texttt{arbitrage\_open}.
\bfield{Rules} Everything is converted to dollars per barrel before it is subtracted; a
crack whose product barrels do not equal its crude barrels is rejected; lots are
rounded to whole contracts and the hedge's lots sum to zero.
\bfield{Acceptance tests} \texttt{code/firm/cracks/tests/}: 7.44 barrels per tonne of
gasoil; the 3-2-1 as a crack; a margin with processing gain; a balanced hedge.
\bfield{Stretch} A crack on futures of different months (crude this month, products
next); a refinery linear programme that chooses the slate from product prices.
\end{build}

\begin{omsources}
\item US Energy Information Administration: \emph{Frequently Asked Questions} (products
from a barrel); \emph{Refinery Capacity Report} and \emph{Today in Energy}, ``U.S.
refining capacity decreased during 2025'' (2026); spot price series via FRED.
\item International Maritime Organization, ``IMO 2020 --- cutting sulphur oxide emissions''.
\item US Environmental Protection Agency, ``Gasoline Reid vapor pressure''.
\item NYMEX Rulebook, chapters 191 (RBOB gasoline) and 150 (NY Harbor ULSD) (Internet Archive copies); ICE Low Sulphur Gasoil
futures specification.
\end{omsources}

\section{Exercises}

\begin{exercise}[$\star$]\label{exo:m3:refined-products-and-cracks:1}
Crude is \qty{75}{\$/\barrel}, gasoline \qty{2.40}{\$/gal}, diesel \qty{2.80}{\$/gal}. Give
the 3-2-1 crack.
\end{exercise}

\begin{exercise}[$\star$]\label{exo:m3:refined-products-and-cracks:2}
Gasoil trades at \qty{900}{\$/t}. What is that per barrel at the density of the ICE
contract?
\end{exercise}

\begin{exercise}[$\star$]\label{exo:m3:refined-products-and-cracks:3}
Why does gasoline trade higher in May than in January, and why does its futures price
step up in the spring?
\end{exercise}

\begin{exercise}[$\star\star$]\label{exo:m3:refined-products-and-cracks:4}
A refinery's yields are 47\% gasoline, 30\% diesel and 30\% other products worth
\qty{60}{\$/\barrel}. With crude at \qty{80}{}, gasoline at \qty{100}{} and diesel at
\qty{110}{\$/\barrel}, what is its gross margin? Why do the yields sum to more than 100\%?
\end{exercise}

\begin{exercise}[$\star\star$]\label{exo:m3:refined-products-and-cracks:5}
New York gasoline is \qty{2.62}{\$/gal}, Gulf gasoline \qty{2.50}{\$/gal}, and moving a
gallon costs \qty{0.08}{\$}. Is the window open? At what Gulf price would it close?
\end{exercise}

\begin{exercise}[$\star\star$]\label{exo:m3:refined-products-and-cracks:6}
How many CL, RB and HO lots lock a 3-2-1 crack on 60{,}000 barrels a day for 30 days?
\end{exercise}

\begin{exercise}[$\star\star\star$]\label{exo:m3:refined-products-and-cracks:7}
\emph{Coding.} With \texttt{crack\_series}, find the month of the highest 3-2-1 crack in
2022 and its value, and the average crack over all months.
\end{exercise}

\begin{exercise}[$\star\star\star$]\label{exo:m3:refined-products-and-cracks:8}
\emph{Find the flaw.} ``Our refinery is hedged: we bought crude futures for all our
throughput.''
\end{exercise}

\section{Problem: Hedging a Refinery's Quarter}

\begin{problem}[{Weekend problem --- locking a record margin}]\label{pb:m3:refined-products-and-cracks:1}
At the end of August 2026 a refinery with 90{,}000 barrels a day of throughput wants
to lock its margin for the next quarter (92 days). Take the August averages as the
futures prices: WTI \qty{83.8976}{\$/\barrel}, gasoline \qty{3.2132}{\$/gal}, diesel
\qty{4.2530}{\$/gal}. Its slate is close to 2 barrels of gasoline for 1 of diesel.

\medskip\noindent\textbf{Part I --- The margin.}
\begin{enumerate}
\item Give the gasoline and diesel prices per barrel.
\item Give the 3-2-1 crack.
\item How does it compare with the average since 2006?
\item Why were distillates so much dearer than gasoline?
\item What was the crack's previous peak, and when?
\end{enumerate}

\medskip\noindent\textbf{Part II --- The hedge.}
\begin{enumerate}[resume]
\item How many barrels of crude will the refinery process in the quarter?
\item How many CL lots does it buy?
\item How many RB and HO lots does it sell?
\item What margin, in dollars, does the hedge lock?
\item If the crack falls to \$25 by delivery, what do the futures gain or lose?
\end{enumerate}

\medskip\noindent\textbf{Part III --- What is not hedged.}
\begin{enumerate}[resume]
\item The refinery's crude is not WTI Cushing. What risk remains?
\item Its products are sold on the Gulf Coast, not in New York. What risk remains?
\item It produces 3 barrels of gasoline for 2 of diesel. How should the hedge change?
\item An unplanned outage stops the refinery for two weeks. What happens to the hedge?
\item What margin must the refinery post on the hedge if the crack widens further?
\end{enumerate}

\medskip\noindent\textbf{Part IV --- Judgement.}
\begin{enumerate}[resume]
\item Why might the refinery hedge only half of the quarter?
\item Who takes the other side of a refiner's crack hedge?
\item Why can a record crack be a good time to hedge and still a hard one?
\item State the \emph{named result}: the lots and the margin the hedge locks.
\item In one sentence: what does a refinery sell?
\end{enumerate}
\end{problem}

\section{Interview questions}

\begin{interviewq}[$\star$ \iqroles{trader}]\label{iq:m3:refined-products-and-cracks:1}
What is a 3-2-1 crack, and why those numbers?
\end{interviewq}

\begin{interviewq}[$\star$ \iqroles{trader, researcher}]\label{iq:m3:refined-products-and-cracks:2}
Gasoil is quoted per tonne and crude per barrel. How do you compare them?
\end{interviewq}

\begin{interviewq}[$\star\star$ \iqroles{researcher}]\label{iq:m3:refined-products-and-cracks:3}
How would you test whether the gasoline crack is seasonal, and trade it if so?
\end{interviewq}

\begin{interviewq}[$\star\star$ \iqroles{trader}]\label{iq:m3:refined-products-and-cracks:4}
Diesel cracks triple in a month while gasoline cracks barely move. What could cause it?
\end{interviewq}

\begin{interviewq}[$\star\star$ \iqroles{risk}]\label{iq:m3:refined-products-and-cracks:5}
A refiner hedges its crack a year ahead. List the risks the hedge does not cover.
\end{interviewq}

\begin{interviewq}[$\star\star\star$ \iqroles{developer, researcher}]\label{iq:m3:refined-products-and-cracks:6}
Design a daily report of refining margins for twenty refineries, each with its own crude
and slate, from market prices.
\end{interviewq}
